\chapter{Feature Composition and Selection}
\label{ch:feature-composition}

This chapter synthesizes Layers~0--4 (Chapters~3--7), adds the remaining layers (calendar, memory, sentiment), and provides the decision framework for which features to use at each forecast horizon.

%% ──────────────────────────────────────────────
\section{Calendar, Memory, and Sentiment (Layers 5--7)}
\label{sec:layers-5-7}
%% ──────────────────────────────────────────────

Layers~5--7 are individually weak but collectively additive.
They are the last 5\% of achievable accuracy.

\subsection*{Layer 5: Calendar Features}

\begin{table}[htbp]
\centering
\caption{Calendar features and when they matter.}
\label{tab:calendar-features}
\small
\begin{tabular}{@{}p{4.8cm}p{8.5cm}@{}}
\toprule
\textbf{Feature} & \textbf{When It Matters} \\
\midrule
FOMC indicator ($-1, 0, +1, +2$ days) &
  Vol compression before announcement, expansion after \\[4pt]
NFP/CPI indicator &
  Macro release effect on intraday and next-day vol \\[4pt]
Options expiry (monthly/quarterly) &
  Pinning and gamma unwind \\[4pt]
Quarter-end rebalancing &
  Last 3 days of quarter; forced portfolio flows \\[4pt]
Earnings proximity &
  Days to next earnings release (single names only) \\[4pt]
Time-of-day &
  U-shape: open and close have highest vol \\[4pt]
Day-of-week &
  Weakening over time but still detectable \\
\bottomrule
\end{tabular}
\end{table}

\subsection*{Layer 6: Memory Features}

\begin{table}[htbp]
\centering
\caption{Memory features capturing long-range dependence and regime state.}
\label{tab:memory-features}
\small
\begin{tabular}{@{}p{4.2cm}p{9cm}@{}}
\toprule
\textbf{Feature} & \textbf{What It Is} \\
\midrule
Fractionally differenced $\RV$ &
  $(1-L)^d \RV$ with $d \approx 0.35$--$0.45$; preserves long memory while stationary \citep{LopezdePrado2018AFML} \\[4pt]
Rolling Hurst exponent &
  $H < 0.15$ = rough/fast mean-reversion; $H > 0.3$ = trending \citep{GJR2018} \\[4pt]
Vol-of-vol &
  $\operatorname{std}(\RV)$ over last 22 days; measures instability of the vol process \\[4pt]
Regime duration &
  Days since last $2\sigma$ $\RV$ spike; acts as a mean-reversion clock \\
\bottomrule
\end{tabular}
\end{table}

\subsection*{Layer 7: Sentiment}

FinBERT news sentiment and negative news count provide 1--3\% $\QLIKE$ improvement, concentrated in crisis periods \citep{AudrinoSignristBallinari2020}.
Include these features only if the data pipeline effort is justified by the target asset class.
For index-level forecasts (SPX, SPY), sentiment adds less value than for single names with idiosyncratic news flow.


%% ──────────────────────────────────────────────
\section{The Diminishing Returns Curve}
\label{sec:diminishing-returns}
%% ──────────────────────────────────────────────

Figure~\ref{fig:diminishing-returns} shows the cumulative accuracy contribution by layer.
The first 20 features (Layers~0--2) achieve 85\% of attainable accuracy.
The remaining 60--100 features add 15\%.

\begin{figure}[htbp]
\centering
\begin{tikzpicture}
\begin{axis}[
  ybar,
  width=\textwidth,
  height=7.5cm,
  bar width=1.4cm,
  ymin=0, ymax=110,
  ylabel={Cumulative accuracy (\%)},
  ylabel style={font=\small},
  symbolic x coords={
    {L0 (5)},
    {+L1 (11)},
    {+L2 (20)},
    {+L3--4 (40)},
    {+L5--7 (80--120)}
  },
  xtick=data,
  xticklabel style={font=\small, align=center},
  yticklabel style={font=\small},
  ytick={0,20,40,60,80,100},
  nodes near coords,
  nodes near coords style={font=\small\bfseries, anchor=south},
  every node near coord/.append style={yshift=1pt},
  enlarge x limits=0.15,
  axis lines=left,
  grid=major,
  grid style={gray!20},
  major tick length=3pt,
]
\addplot[fill=defblue!70, draw=defblue] coordinates {
  ({L0 (5)}, 55)
  ({+L1 (11)}, 70)
  ({+L2 (20)}, 85)
  ({+L3--4 (40)}, 95)
  ({+L5--7 (80--120)}, 100)
};
\draw[dashed, warnred, thick] ({axis cs:{+L2 (20)},85}) -- ({axis cs:{+L5--7 (80--120)},85})
  node[midway, below=2pt, font=\footnotesize, text=warnred, fill=white, inner sep=1pt] {85\% with 20 features};
\end{axis}
\end{tikzpicture}
\caption{Cumulative accuracy by feature layer. Parentheses show total feature count at each stage. The dashed line marks the 85\% threshold achievable with Layers~0--2 and a Ridge regression.}
\label{fig:diminishing-returns}
\end{figure}


%% ──────────────────────────────────────────────
\section{Feature Priority by Forecast Horizon}
\label{sec:horizon-priority}
%% ──────────────────────────────────────────────

\begin{table}[htbp]
\centering
\caption{Dominant features and ML value-add by forecast horizon.}
\label{tab:horizon-priority}
\small
\begin{tabular}{@{}p{2.8cm}p{4.8cm}p{5.8cm}@{}}
\toprule
\textbf{Horizon} & \textbf{Dominant Features} & \textbf{Where ML Adds Value} \\
\midrule
Intraday (10min--1hr) &
  Microstructure (L3): price acceleration, OBI, spread &
  Trees with 600+ features \\[6pt]
1 day &
  HAR core (L0) + RQ + asymmetry (L1) &
  HARQ nearly optimal; ML adds ${\sim}5\%$ \\[6pt]
1 week &
  Options (L2) + cross-asset (L4) &
  VRP + skew: 5--10\% over pure RV models \\[6pt]
1 month &
  VRP (L2) + macro (L4) + Hurst (L6) &
  Options have max informational advantage \\
\bottomrule
\end{tabular}
\end{table}

\begin{warning}[The 1-Day Horizon Is the Hardest to Beat]
At the 1-day horizon, HARQ with 5 features is nearly optimal.
ML models gain ${\sim}5\%$ $\QLIKE$ improvement at best, and much of that comes from the RQ interaction that HARQ already captures.
The largest ML gains are at intraday and weekly+ horizons where linear models lack the capacity to exploit high-dimensional feature sets.
\end{warning}


%% ──────────────────────────────────────────────
\section{Feature Engineering Principles}
\label{sec:engineering-principles}
%% ──────────────────────────────────────────────

\begin{itemize}[leftmargin=1.5em, itemsep=6pt]
  \item \textbf{Triple expansion.}
    For each base quantity, compute \{level, change, $z$-score\} systematically.
    This triples the feature count and captures state, direction, and unusualness in a single pass.

  \item \textbf{Horizon-dependent selection.}
    Drop microstructure features at the monthly horizon (noise at that scale).
    Drop calendar features at intraday (FOMC day is already known by 9:30am).
    Match the feature set to the prediction target.

  \item \textbf{Trees handle redundancy naturally.}
    Correlated features split across nodes without the multicollinearity problems that plague linear models.
    No need to pre-decorrelate or drop correlated pairs before training a gradient-boosted ensemble.

  \item \textbf{Feature importance from a single model is unstable.}
    When features are redundant, a single LightGBM run produces noisy importance rankings that shift across random seeds.
    Rashomon analysis (Chapter~12) addresses this by examining the set of near-optimal models rather than one point estimate.
\end{itemize}


%% ──────────────────────────────────────────────
%% Boxes
%% ──────────────────────────────────────────────

\begin{keyidea}[Features Over Models]
The feature set matters more than the model.
Layers~0--2 achieve 85\% of attainable accuracy with 20 features and a Ridge regression.
The remaining 15\% requires 60--100 more features and careful model architecture.
Invest in feature engineering first; optimize the model second.
\end{keyidea}
